%% paper/sections/05_metronome_algorithm.tex
%% §5 METRONOME 알고리즘 — 5-stage pseudocode + Theorem 1, 2

\section{\algname 알고리즘}
\label{sec:metronome}

본 절은 본 논문의 핵심 contribution인 \algname 알고리즘을 형식화한다.
\algname 은 \textbf{M}ulti-domain \textbf{E}mpirically-derived \textbf{T}iming-aligned \textbf{R}esource \textbf{O}rchestration with \textbf{N}egative-result-driven \textbf{O}ptimization for \textbf{M}odular \textbf{E}xecution의 약자이며, 5단계 파이프라인(\textsc{Tempo} $\rightarrow$ \textsc{Chord} $\rightarrow$ \textsc{Meter} $\rightarrow$ \textsc{Accent} $\rightarrow$ \textsc{Ritardando})으로 구성된다.
음악적 명명은 GPU step 주기를 박자(\emph{tempo})로 간주하고 CPU 보조 작업을 그 박자에 정렬하는 핵심 directive를 직관적으로 묶는다.

%% figures/fig1_metronome_pipeline.tex
%% Fig.1 — METRONOME 5-stage pipeline flow diagram (TikZ)
%% \input from §5.

\begin{figure}[t]
\centering
\begin{tikzpicture}[
    node distance=0.42cm and 0.50cm,
    stage/.style={rectangle, draw, rounded corners=2pt, minimum width=2.4cm,
                  minimum height=0.65cm, align=center, font=\small},
    decision/.style={diamond, draw, aspect=2.4, minimum width=2.4cm,
                     minimum height=0.4cm, align=center, font=\footnotesize},
    annot/.style={font=\footnotesize\itshape, text=gray!70!black, align=left},
    arrow/.style={-{Stealth[length=2.0mm]}, thick},
]
% Stages — top to bottom
\node[stage, fill=blue!8] (s1) {Stage 1 — \textsc{Tempo}\\\footnotesize $\tgpu \gets$ probe step period};
\node[stage, fill=blue!8, below=of s1] (s2) {Stage 2 — \textsc{Chord}\\\footnotesize cross-domain pair};
\node[stage, fill=blue!8, below=of s2] (s3) {Stage 3 — \textsc{Meter}\\\footnotesize $\tcpu \gets k \cdot \tgpu$, $k \geq 2$};
\node[stage, fill=blue!8, below=of s3] (s4) {Stage 4 — \textsc{Accent}\\\footnotesize work pattern: branchy/regular};
\node[stage, fill=green!12, below=of s4] (s5) {Stage 5 — \textsc{Ritardando}\\\footnotesize continuous adaptive};
\node[decision, fill=yellow!12, below=of s5] (drift) {drift $> 20\%$?};

% Arrows
\draw[arrow] (s1) -- (s2);
\draw[arrow] (s2) -- (s3);
\draw[arrow] (s3) -- (s4);
\draw[arrow] (s4) -- (s5);
\draw[arrow] (s5) -- (drift);
\draw[arrow] (drift.east) -- ++(0.7,0) |- node[pos=0.25, right, font=\footnotesize] {yes (a tempo)} (s1.east);
\draw[arrow] (drift.west) -- ++(-0.7,0) |- node[pos=0.25, left, font=\footnotesize] {no} (s5.west);

% Annotations on the right
\node[annot, right=0.20cm of s1] {박자 측정};
\node[annot, right=0.20cm of s2] {도메인 분리};
\node[annot, right=0.20cm of s3] {주기 정렬};
\node[annot, right=0.20cm of s4] {강세 결정};
\node[annot, right=0.20cm of s5] {지속 적응};

\end{tikzpicture}
\caption{\algname 알고리즘의 5단계 파이프라인. 음악적 명명은 각 단계의 역할을 직관적으로 묶는다: \textsc{Tempo}는 박자 측정, \textsc{Chord}는 교차 도메인 기법 쌍 선택, \textsc{Meter}는 주기 정렬, \textsc{Accent}는 작업 패턴 강세 결정, \textsc{Ritardando}는 runtime drift 시 tempo 재측정이다.}
\label{fig:metronome-pipeline}
\end{figure}


\subsection{단계별 설명}
\label{subsec:stages}

\subsubsection{Stage 1 --- \textsc{Tempo} (박자 측정)}

GPU step 주기 $\tgpu$를 측정한다.
초기 $n$개(본 연구는 $n=10$)의 step에 대해 wall-clock을 측정하여 평균값을 계산한다.

\begin{equation}
\tgpu = \frac{1}{n} \sum_{i=1}^{n} t_{\text{step},i}
\end{equation}

이 값은 \textsc{Meter}와 \textsc{Ritardando} 단계의 기준이 된다.

\subsubsection{Stage 2 --- \textsc{Chord} (교차 도메인 기법 쌍)}

CPU 최적화 기법 후보 집합 $\mathcal{L}$에서 \emph{서로 다른 도메인}에 속하는 두 기법 $L_{\text{primary}}, L_{\text{secondary}}$를 선택한다.
본 연구에서 도메인은 다음과 같이 분류된다.

\begin{itemize}
  \item $\mathrm{system}$ : OS-level 격리 (NUMA pin, cgroup, hugepages, taskset)
  \item $\mathrm{work}$ : CPU 보조 작업 (regular/branchy 패턴, period $\tcpu$)
  \item $\mathrm{kernel}$ : 직접 커널 대체 (AVX-512, AMX) --- 본 연구의 측정 결과 critical path 침범으로 일관 net-negative이므로 \algname 에서 제외
  \item $\mathrm{scheduling}$ : sub-step CPU 스케줄링 --- 자원 도메인 간섭으로 net-negative이므로 제외
\end{itemize}

선택 제약은 $\dom(L_{\text{primary}}) \neq \dom(L_{\text{secondary}})$이며, 이는 \S\ref{subsec:theorem1}의 Theorem~\ref{thm:superposition}에 의해 정당화된다.
본 연구의 최적 구성은 $L_{\text{primary}} = $ NUMA pin (system), $L_{\text{secondary}} = $ branchy auxiliary work (work)이다.

\subsubsection{Stage 3 --- \textsc{Meter} (주기 정렬)}

CPU 보조 작업의 주기 $\tcpu$를 $\tgpu$의 정수배로 정렬한다.

\begin{equation}
\tcpu = \max(k \cdot \tgpu, \tcpu^{\min}) \quad \text{where } k \in \{2, 3\}
\label{eq:meter}
\end{equation}

본 연구의 측정 환경에서 $\tgpu \approx 40$\,ms, $\tcpu^{\min} = 100$\,ms로 설정하여 $k \approx 2.5$가 된다.
이 정렬은 \S\ref{subsec:theorem2}의 Theorem~\ref{thm:alignment}에 의해 정당화된다.

%% figures/fig2_step_alignment.tex
%% Fig.2 — Step boundary timing diagram (단일 컬럼 fit)

\begin{figure}[t]
\centering
\resizebox{0.95\columnwidth}{!}{%
\begin{tikzpicture}[
    x=0.40cm, y=0.55cm,
    gpu/.style={rectangle, draw, fill=blue!20, minimum height=0.42cm,
                font=\scriptsize, inner sep=1pt},
    cpu/.style={rectangle, draw, fill=green!22, minimum height=0.38cm,
                font=\scriptsize, inner sep=1pt},
    axis/.style={-{Stealth[length=1.6mm]}},
    annot/.style={font=\scriptsize, text=gray!60!black},
]
% Time axis
\draw[axis] (-0.2, 0) -- (16, 0) node[right, font=\scriptsize] {ms};
\foreach \x/\t in {0/0, 4/40, 8/80, 12/120, 16/160} {
    \draw (\x, 0) -- (\x, -0.1);
    \node[below, font=\tiny] at (\x, -0.1) {\t};
}

% GPU step lane (label on left)
\node[font=\scriptsize, anchor=east] at (-0.2, 0.7) {GPU};
\foreach \i in {0,1,2,3} {
    \pgfmathsetmacro{\xs}{\i*4}
    \node[gpu, anchor=south west, minimum width={4*0.40cm-1pt}]
        at (\xs, 0.4) {\(\tgpu\)};
}

% CPU aux lane
\node[font=\scriptsize, anchor=east] at (-0.2, 1.7) {CPU};
\node[cpu, anchor=south west, minimum width={3*0.40cm-1pt}] at (0, 1.4) {aux};
\node[cpu, anchor=south west, minimum width={3*0.40cm-1pt}] at (10, 1.4) {aux};

% Cycle marker
\draw[<->, thick, gray!70!black] (0, 2.4) -- (10, 2.4)
    node[midway, above=1pt, annot] {$\tcpu = k\tgpu = 100$\,ms ($k\!=\!2.5$)};
\draw[<->, thick, gray!70!black] (0, -0.6) -- (4, -0.6)
    node[midway, below=1pt, annot] {$\tgpu \approx 40$\,ms};

% Boundary alignment
\foreach \x in {0, 4, 8, 12} {
    \draw[dotted, gray!50] (\x, -0.6) -- (\x, 2.4);
}
\end{tikzpicture}%
}
\caption{\textsc{Meter} 단계의 박자 정렬 (Theorem~\ref{thm:alignment}). GPU step 주기 $\tgpu \approx 40$\,ms 에 대해 CPU 보조 작업 주기 $\tcpu = 100$\,ms ($k=2.5$) 가 step boundary 와 정수배로 정렬되어 critical path 침범 없이 amortize window 를 확보한다.}
\label{fig:step-alignment}
\end{figure}


\subsubsection{Stage 4 --- \textsc{Accent} (작업 패턴 강세)}

CPU 보조 작업의 \emph{패턴}을 결정한다.
GPU의 PCIe pressure $\rho_{\text{pcie}}$를 측정하여 임계값 $\tau$와 비교한다.

\begin{equation}
\text{pattern} = \begin{cases}
\mathrm{BRANCHY} & \text{if } \rho_{\text{pcie}} > \tau \\
\mathrm{REGULAR} & \text{otherwise}
\end{cases}
\end{equation}

PCIe pressure가 높은 환경(GPU 8장 + 텐서 병렬)에서는 branchy 패턴(rank/sort 등 분기 예측 불가 코드)이 CPU prefetcher의 적극 활성을 억제하여 GPU PCIe traffic과의 contention을 회피한다.
이 경험적 발견은 직관(고빈도 regular work가 더 효과적일 것)과 반대이며, \S\ref{sec:mechanism}에서 메커니즘을 자세히 분석한다.

\subsubsection{Stage 5 --- \textsc{Ritardando} (지속 적응)}

런타임 동안 step 주기 $\tgpu$의 drift를 모니터링한다.

\begin{equation}
\delta(t) = \frac{|t_{\text{step}}(t) - \tgpu|}{\tgpu}
\end{equation}

$\delta(t) > 0.2$이면 model 또는 workload가 변경된 것으로 간주하여 Stage 1로 복귀(\emph{a tempo}).
이 단계는 model swap, 다른 워크로드 mix 등 변화하는 환경에서 \algname 의 강건성을 보장한다.

\subsection{Theorem 1: 교차 도메인 중첩 (Cross-Domain Superposition)}
\label{subsec:theorem1}

\begin{theorem}[Cross-Domain Superposition]
\label{thm:superposition}
두 최적화 기법 $L_A, L_B$의 stacking 효과는 다음과 같이 분해된다.

\begin{equation}
\Delta(L_A \oplus L_B) = \Delta(L_A) + \Delta(L_B) - \varepsilon(A, B)
\label{eq:superposition}
\end{equation}

여기서 \emph{도메인 간섭 비용} $\varepsilon(A, B)$는 도메인 관계에 따라 다음을 만족한다.

\begin{equation}
\varepsilon(A, B) = \begin{cases}
\varepsilon_{\text{cross}} & \text{if } \dom(A) \neq \dom(B) \\
\varepsilon_{\text{same}} & \text{if } \dom(A) = \dom(B)
\end{cases}
\end{equation}

본 연구의 경험적 측정에 따라

\begin{equation}
\varepsilon_{\text{cross}} \approx 0.55\,\mathrm{pp} \ll \varepsilon_{\text{same}} \approx 1.18\,\mathrm{pp}
\end{equation}

이다. 또한 동일 도메인 3-기법 stacking은 초선형 destructive interference로 확장되어 $\varepsilon_{\text{same}}^{(3)} \approx 4.56\,\mathrm{pp}$의 비용을 발생시킨다.
\end{theorem}

\begin{proof}[증명 개요 (경험적)]
본 정리는 \S\ref{sec:results}의 Table~\ref{tbl:superposition}의 측정 결과에서 도출된다.
구체적으로 동일 도메인 OS isolation 4-기법 stack에서 측정된 residual은 단일 기법 합 대비 $-1.18$\,pp, 교차 도메인 NUMA + auxiliary work pair에서는 $-0.55$\,pp가 측정되었다.
이론적 근거는 동일 도메인 기법이 같은 OS 자원(context-switch slot, cache)을 경쟁하기 때문으로 해석된다.
\end{proof}

\subsection{Theorem 2: 단계 경계 정렬 (Step-Boundary Alignment)}
\label{subsec:theorem2}

\begin{theorem}[Step-Boundary Alignment]
\label{thm:alignment}
CPU 보조 작업 주기 $\tcpu$를 GPU step 주기 $\tgpu$의 정수배 $k \geq 2$로 정렬할 때 amortize window가 충분히 커져 다음의 정렬 이득이 발생한다.

\begin{equation}
\Delta_{\text{align}}(\tcpu) \propto \min\!\left(\frac{\tcpu}{\tgpu} - 1, \, k_{\max} - 1\right)
\label{eq:alignment}
\end{equation}

여기서 $k_{\max}$는 model의 보조 작업 amortize 가능 최대 배수이며, 본 연구의 측정에서 $k_{\max} \approx 3$이다.
\end{theorem}

\begin{proof}[증명 개요]
$\tcpu \ll \tgpu$이면 보조 작업이 step 내부에 끼어들어 critical path 침범으로 net-negative이다.
$\tcpu \approx \tgpu$이면 step boundary 정렬이 정확하지만 amortize 효율이 낮다.
$\tcpu = k \cdot \tgpu$, $k \geq 2$이면 idle window가 충분히 커져 보조 작업이 완료될 수 있으며, $k$가 커질수록 step당 amortize 비율 $(k-1)/k$가 1에 수렴한다.
단 $k > k_{\max}$이면 idle window 대부분이 낭비되어 추가 이득이 없다.

본 연구의 측정에서 $\tgpu = 40$\,ms, $\tcpu = 100$\,ms ($k=2.5$)일 때 $\Delta_{\text{align}} = +5.28\,\%$로 최대 효과를 나타냈다.
$\tcpu = 10$\,ms ($k = 0.25$)에서 $+0.98\,\%$, $\tcpu = 20$\,ms ($k = 0.5$)에서 1-run $+1.66\,\%$에 그쳤다.
\end{proof}

\subsection{알고리즘 의사 코드}

\algname 은 두 함수로 분해된다.
\textsc{Initialize} 는 박자 측정과 화음/주기/강세 결정의 \emph{초기화 단계}이며 (Algorithm~\ref{alg:metronome-init}), \textsc{Serve} 는 runtime 실행과 drift 적응의 \emph{실행 단계}이다 (Algorithm~\ref{alg:metronome-serve}).
두 함수의 분리는 hyperparameter 도출이 한 번 일어나고 runtime 은 그 결과를 안정적으로 활용함을 명확히 한다.

\begin{algorithm}[t]
\caption{\algname --- Initialization phase.}
\label{alg:metronome-init}
\begin{algorithmic}[1]
\Function{Initialize}{$M$, $N$, $\tau$}
  \LineComment{$M$: model config, $N$: NUMA topology, $\tau$: PCIe threshold}
  \LineComment{\textsc{Stage 1 — Tempo}: 박자 측정}
  \State $S \gets \emptyset$
  \For{$i = 1$ \textbf{to} $10$}
    \State $S \gets S \cup \{$\Call{ProbeStepLatency}{$M$}$\}$
  \EndFor
  \State $\tgpu \gets \frac{1}{|S|}\sum_{s \in S} s$
  \LineComment{\textsc{Stage 2 — Chord}: 교차 도메인 화음}
  \State $L_p \gets$ \Call{NumaPin}{$N$} \Comment{$\dom(L_p) = \mathrm{system}$}
  \State $L_s \gets$ \Call{AuxiliaryWorkHandle}{$\;$} \Comment{$\dom(L_s) = \mathrm{work}$}
  \If{$\dom(L_p) = \dom(L_s)$}
    \State \Return \textsc{Error} \Comment{Theorem~\ref{thm:superposition} 위반}
  \EndIf
  \LineComment{\textsc{Stage 3 — Meter}: 주기 정렬}
  \State $\tcpu \gets \max(2 \cdot \tgpu, \, 100\,\mathrm{ms})$
  \Comment{$k = \lceil \tcpu / \tgpu \rceil \in \{2,3\}$}
  \LineComment{\textsc{Stage 4 — Accent}: 작업 패턴 결정}
  \State $\rho \gets$ \Call{MeasurePcieUtilization}{$M$}
  \If{$\rho > \tau$}
    \State $P \gets \mathrm{BRANCHY}$
  \Else
    \State $P \gets \mathrm{REGULAR}$
  \EndIf
  \State \Return $(\tgpu, L_p, L_s, \tcpu, P)$
\EndFunction
\end{algorithmic}
\end{algorithm}

\begin{algorithm}[t]
\caption{\algname --- Serve phase with adaptive re-probing.}
\label{alg:metronome-serve}
\begin{algorithmic}[1]
\Function{Serve}{$M$, $N$, $\tau$, $\delta_{\max}$}
  \LineComment{$\delta_{\max}$: drift threshold (default 0.20)}
  \State $(\tgpu, L_p, L_s, \tcpu, P) \gets$ \Call{Initialize}{$M$, $N$, $\tau$}
  \State \Call{Deploy}{$L_p, L_s$}
  \State $h \gets$ \Call{LaunchAuxiliaryWorker}{$L_s, P, \tcpu$}
  \LineComment{\textsc{Stage 5 — Ritardando}: 지속 적응}
  \While{server is running}
    \State $t \gets$ \Call{ObserveCurrentStepLatency}{$M$}
    \State $\delta \gets |t - \tgpu| / \tgpu$
    \If{$\delta > \delta_{\max}$}
      \State \Call{Stop}{$h$} \Comment{a tempo: 박자 재측정}
      \State $(\tgpu, L_p, L_s, \tcpu, P) \gets$ \Call{Initialize}{$M$, $N$, $\tau$}
      \State $h \gets$ \Call{LaunchAuxiliaryWorker}{$L_s, P, \tcpu$}
    \EndIf
    \State \Call{Sleep}{$\tgpu$} \Comment{step boundary 마다 1회 점검}
  \EndWhile
\EndFunction
\end{algorithmic}
\end{algorithm}

알고리즘~\ref{alg:metronome-init} 의 \textsc{Initialize} 는 한 번 호출되어 $\tgpu, L_p, L_s, \tcpu, P$ 의 5 가지 hyperparameter 를 반환한다.
이 중 \textsc{Chord} 단계 (line 8--11) 는 Theorem~\ref{thm:superposition} 의 $\dom(L_p) \neq \dom(L_s)$ 제약을 강제로 검사한다.
\textsc{Meter} 단계 (line 12) 는 Theorem~\ref{thm:alignment} 의 정수배 정렬을 보장하기 위해 $k \geq 2$ 의 lower bound 와 $\tcpu^{\min} = 100$\,ms 의 floor 를 동시에 적용한다.

알고리즘~\ref{alg:metronome-serve} 의 \textsc{Serve} 는 runtime adaptive loop 이다.
매 $\tgpu$ 주기마다 step latency 의 drift 를 점검하여 $\delta > \delta_{\max}$ 일 때 박자 재측정을 trigger 한다.
이는 model swap, workload mix 변화, 또는 하드웨어 thermal throttle 같은 environmental change 에 대한 강건성을 제공한다.
